\documentclass[10pt,a4paper]{amsart}
\usepackage[margin=19mm]{geometry}
\usepackage{amsmath,amssymb,mathtools}
\usepackage[scr=rsfs,cal=euler]{mathalfa}
\usepackage{xfrac}
\usepackage[unicode,hidelinks,bookmarks=false]{hyperref}
\newcommand{\ie}{i.e.\ }
\newcommand{\cf}{cf.\ }
\newcommand{\resp}{respectively\ }
\newcommand{\Subsection}{\subsection}
\providecommand{\nobreakdash}{\nobreak\hbox{-}}
\setlength{\emergencystretch}{2em}
\multlinegap0pt
\usepackage{bm}
\usepackage{bbm}
\usepackage{dsfont}
\usepackage{xspace}
\usepackage{cancel}
\usepackage{mathtools}
\usepackage{amsthm}
\makeatletter
\@ifundefined{theorem}{\newtheorem{theorem}{Theorem}}{}
\@ifundefined{lemma}{\newtheorem{lemma}[theorem]{Lemma}}{}
\@ifundefined{proposition}{\newtheorem{proposition}[theorem]{Proposition}}{}
\makeatother
\newcommand{\C}{\mathbb{C}}
\newcommand{\N}{\mathbb{N}}
\newcommand{\ODT}[1][]{   ^{\widetilde{\mathrm{OD}}#1}   }
\newcommand{\OD}[1][]{   ^{\mathrm{OD}#1}   }
\newcommand{\R}{\mathbb{R}}
\newcommand{\Z}{\mathbb{Z}}
\newcommand{\dd}{\mathrm{d}}
\newcommand{\e}{\mathrm{e}}
\newcommand{\fd}{r}
\newcommand{\mL}{\mathcal{L}}
\newcommand{\tr}{\mathrm{tr}}
\newcommand{\w}{\omega}
\renewcommand{\i}{\mathrm{i}}
\title[The Colored Hofstadter Butterfly as a Many-Body Quantum Hall]{The Colored Hofstadter Butterfly as a Many-Body Quantum Hall Phase Diagram}
\author{Giovanna Marcelli, Tadahiro Miyao, Domenico Monaco, Stefan Teufel, Marius Wesle}
\date{Source: arXiv:2606.26256v1 (CC BY 4.0)}
\begin{document}
\maketitle

\noindent\emph{Source and licence.} The Colored Hofstadter Butterfly as a Many-Body Quantum Hall Phase Diagram. Version arXiv:2606.26256v1 (CC BY 4.0), distributed under the Creative Commons Attribution 4.0 International licence, as recorded in the arXiv licence metadata for this exact version and shown on the abstract page https://arxiv.org/abs/2606.26256v1. Full rights notice: licenses/arxiv-2606.26256-CC-BY-4.0.txt.\par\smallskip

\noindent\textit{Editorial scope.} Counted unit: the Lemma whose source label is \texttt{lem:diagonal-block-products-trace-class} in the article (source0.tex lines 1397--1411, source label \texttt{lem:diagonal-block-products-trace-class}), delivered together with the article's own complete original proof of it (source0.tex lines 1413--1479, one \texttt{proof} environment of 67 lines). No mathematics is written, repaired or re-derived: statement and proof are the article's own text line for line. Required context retained uncounted: prop:dcformulas (lines 1086--1106); lambdajOD (lines 1157--1159); lem:diagonal-block-strip-local (lines 1204--1218). Every definition, display and label of those lines is retained unchanged. Omitted from this package and not redistributed: 1--1085, 1107--1156, 1160--1203, 1219--1396, 1412--1412, 1480--2300. Nothing inside a retained excerpt is omitted or altered: every retained body is delivered exactly as the article's lines are written, so no omission receipt of any kind accompanies this package. Citations: the retained excerpts carry no citation command at all, so no bibliography entry of the article is transcribed and no reference is redistributed. Reference adaptation: none was needed and none was made; every reference inside the delivered unit resolves inside it, so no unresolved reference or citation is produced. Printed slips kept literal: no printed word, symbol, hypothesis or number of the counted statement or of its proof was altered. Layout and class: the article's own class is \texttt{article} with packages including xcolor, inputenc, babel, amsmath, amssymb, amsthm, comment, hyperref, multicol, tikz, tocbibind, graphicx, bbold, csquotes; the wrapper is the lane's pinned \texttt{amsart} editorial wrapper at 10pt on A4, which restates the theorem-like environments the retained text needs and adds the article's own macro definitions that the retained text needs (\texttt{\textbackslash C}, \texttt{\textbackslash N}, \texttt{\textbackslash OD}, \texttt{\textbackslash ODT}, \texttt{\textbackslash R}, \texttt{\textbackslash Z}, \texttt{\textbackslash dd}, \texttt{\textbackslash e}, \texttt{\textbackslash fd}, \texttt{\textbackslash i}, \texttt{\textbackslash mL}, \texttt{\textbackslash tr}, \texttt{\textbackslash w}), with extra wrapper packages (bm, bbm, dsfont, xspace, cancel, mathtools). The delivered numbering is the unit's own. Assets: no retained span contains a figure environment, a graphics-inclusion command, a PDF-page inclusion, a table or a TikZ picture; no image file is copied into this package, the package declares no asset dependency and no image file is redistributed with it. Licence: version arXiv:2606.26256v1 of this article is distributed under the Creative Commons Attribution 4.0 International licence, as recorded in the arXiv licence metadata for this exact version and shown on the abstract page https://arxiv.org/abs/2606.26256v1. A full rights notice is delivered at licenses/arxiv-2606.26256-CC-BY-4.0.txt. Characters: every retained excerpt is pure ASCII, so the delivered engine, XeLaTeX with its default Latin Modern text font, reports no missing character.
\begin{proposition}\label{prop:dcformulas}
 Let $h\in\mL(\ell^2(\Z^2,\C^\fd))$ be self-adjoint  and $\mu\in\R\setminus\sigma(h)$ such that $g:= \mathrm{dist}(\mu,\sigma(h))>0$.  Let $P:= \chi_{(-\infty,\mu)}(h)$. Assume that the kernel $h(x,y)$ of $h$ has uniform exponential decay, i.e.,
 \begin{equation}\label{hexp}
 \exists c,C>0\;\forall x,y\in\Z^2\quad \|h(x,y)\|\leq C\e^{-c\|x-y\|}\,.
 \end{equation}

\begin{enumerate}
    \item[(i)] 
 Then $H:=\dd\Gamma(h-\mu)\in B_{\exp}$ and the unique quasi-free state $\w$ with $\w(a^*_{x,j}a_{y,i}) = P(x,y)_{ji}$   for all $x,y\in\Z^2$ and $j,i\in\{1,\ldots,\fd\}$ is the unique ground state of $\mL_{H}$ and has a gap of size at least $g$.
\item[(ii)] 
Let $\lambda_j(x):= \chi_{x_j\geq 0}(x)$, $j\in\{1,2\}$,  denote the one-body switch-functions, which define  corresponding multiplication operators $\lambda_j$ on $\ell^2(\Z^2,\C^\fd)$, and $\Lambda_j := \dd\Gamma(\lambda_j)$ their many-body versions. Then
 \[
\i\,\w([\Lambda_2\OD,\Lambda_1\OD]) = \i\,\tr (P [\lambda_2\ODT,\lambda_1\ODT]) \in \tfrac{1}{2\pi} \Z\,,
 \]
 where $\lambda_j\ODT := P\lambda_jP^\perp +P^\perp\lambda_jP $.
\item[(iii)] 
 If $H$ is periodic with respect to some translation $T$, then also
 \[ \overline\w([X_2\OD,X_1\OD])=\w([\Lambda_2\OD,\Lambda_1\OD])\,.
 \]
 \end{enumerate}
\end{proposition}

\begin{equation} \label{lambdajOD}
\lambda_j\OD :=  \i \,  \int_{\R}\dd s \, W_g(s)\, \e^{\i s h} \,  [\lambda_j, h]\e^{-\i s h}\,.
\end{equation}

\begin{lemma}\label{lem:diagonal-block-strip-local}
Let $h$, $P$, and $\lambda_j$ be as in Proposition~\ref{prop:dcformulas}
and define  $\lambda_j\OD$ as in~\eqref{lambdajOD}. 
Then, for every $N\in\N$, there exists $C_N<\infty$ such that
\[
    \bigl\|(P\lambda_j\OD P)(x,y)\bigr\|
    \le
    C_N
    \langle \|x-y\|\rangle^{-N}
    \langle |x_j|\rangle^{-N}
    \langle |y_j|\rangle^{-N}
\]
for all $x,y\in\Z^2$. Here $\langle r\rangle := 1+r$, and
$\|\cdot\|$ denotes the $\ell^\infty$-norm on $\Z^2$.
\end{lemma}

\begin{lemma}\label{lem:diagonal-block-products-trace-class}
With the notation of Lemma~\ref{lem:diagonal-block-strip-local}, the operator
\[
    (P\lambda_1\OD P)(P\lambda_2\OD P)
    =
    P\lambda_1\OD P\lambda_2\OD P
\]
and thus also its adjoint
\[
    (P\lambda_2\OD P)(P\lambda_1\OD P)
    =
    P\lambda_2\OD P\lambda_1\OD P
\]
is trace class.
\end{lemma}

\begin{proof}
For $j=1,2$ set
    $D_j:=P\lambda_j\OD P$.
By Lemma~\ref{lem:diagonal-block-strip-local}, for every $N\in\N$ there is
$C_N<\infty$ such that
\[
    \|D_1(x,z)\|
    \le
    C_N
    \langle \|x-z\|\rangle^{-N}
    \langle |x_1|\rangle^{-N}
    \langle |z_1|\rangle^{-N}
\]
and
\[ 
    \|D_2(z,y)\|
    \le
    C_N
    \langle \|z-y\|\rangle^{-N}
    \langle |z_2|\rangle^{-N}
    \langle |y_2|\rangle^{-N}\,.
\]
Choose $N>2$. Then
\begin{align*}
    \hspace{2em}&\hspace{-2em}
    \sum_{x,y\in\Z^2}\|(D_1D_2)(x,y)\|
    \\
    &\leq
    C_N^2
    \sum_{x,y,z\in\Z^2}
    \langle \|x-z\|\rangle^{-N}
    \langle |x_1|\rangle^{-N}
    \langle |z_1|\rangle^{-N}    
    \langle \|z-y\|\rangle^{-N}
    \langle |z_2|\rangle^{-N}
    \langle |y_2|\rangle^{-N}\,.
\end{align*}
For fixed $z$, since $N>2$,
\[
    \sum_{x\in\Z^2}
    \langle \|x-z\|\rangle^{-N}
    \langle |x_1|\rangle^{-N}
    \le
    \sum_{x\in\Z^2}
    \langle \|x-z\|\rangle^{-N}
    \le C_N'\,,
\]
uniformly in $z$. Similarly,
\[
    \sum_{y\in\Z^2}
    \langle \|z-y\|\rangle^{-N}
    \langle |y_2|\rangle^{-N}
    \le C_N'\,,
\]
uniformly in $z$. Therefore
\[
    \sum_{x,y\in\Z^2}\|(D_1D_2)(x,y)\|
    \leq
    (C_N')^2
    \sum_{z\in\Z^2}
    \langle |z_1|\rangle^{-N}
    \langle |z_2|\rangle^{-N}
    <\infty\,.
\]
Thus $D_1D_2$ has an absolutely summable matrix kernel. Since the fiber
dimension is finite, this implies that $D_1D_2$ is trace class. 
\end{proof}

\end{document}
